%% Part 2: Problem Formulation + Method
%% Sections 3–4 of the NeurIPS 2026 positioning paper

% ============================================================
\section{Problem Formulation}
\label{sec:formulation}
% ============================================================

\paragraph{Frozen-Model Translation Setting.}
Let $\mathcal{X}_S \sim P_S$ denote source-domain clinical time series (eICU) and $\mathcal{X}_T \sim P_T$ denote the target domain (MIMIC-IV).
A target-domain predictor $f_T : \mathbb{R}^{T \times F} \to \mathbb{R}^{T \times C}$ has been trained on $\mathcal{X}_T$ and is \emph{entirely frozen}: all parameters satisfy $\nabla_\phi f_T = 0$ and are never updated.
We learn a translator $g_\theta : \mathbb{R}^{T \times F} \to \mathbb{R}^{T \times F}$ that maps source inputs into the target input space so that the composite system $f_T \circ g_\theta$ performs well on source labels $y_S$.
The composite training objective is:
\begin{equation}
\label{eq:composite}
\mathcal{L}(\theta) \;=\; \mathcal{L}_{\mathrm{task}}\!\bigl(f_T(g_\theta(x_S)),\, y_S\bigr) \;+\; \lambda_{\mathrm{fid}}\,\mathcal{L}_{\mathrm{fid}}\!\bigl(g_\theta(x_S),\, x_S\bigr) \;+\; \lambda_{\mathrm{range}}\,\mathcal{L}_{\mathrm{range}}\!\bigl(g_\theta(x_S)\bigr),
\end{equation}
where $\mathcal{L}_{\mathrm{task}}$ is cross-entropy (classification) or MSE (regression) evaluated through the frozen predictor, $\mathcal{L}_{\mathrm{fid}}$ is an MSE fidelity penalty that regularises the translator toward identity, and $\mathcal{L}_{\mathrm{range}}$ penalises outputs outside the physiological range observed in $\mathcal{X}_T$.

The frozen constraint is motivated by regulatory and deployment considerations.
Under the FDA Software-as-a-Medical-Device (SaMD) ``locked algorithm'' pathway, the downstream predictor is a validated, auditable artifact whose parameters must not change post-deployment~\cite{fdaaiml2021}.
By learning only the input mapping $g_\theta$, clinical sites can adapt a certified predictor to local data distributions without re-validation of $f_T$.
This setting is strictly more constrained than test-time adaptation methods that update batch-normalisation statistics~\cite{wang2021tent} (which modify running mean and variance, altering the predictor's output for identical inputs) or partial model parameters~\cite{liang2020shot} (which retrain the feature extractor while freezing only the classifier head), and distinct from adapter-based approaches~\cite{houlsby2019adapters} that insert learnable modules into the predictor's computation graph. In our setting, the predictor is treated as a fixed, deterministic function---no internal parameter, statistic, or architectural element is modified.

\paragraph{Gradient Alignment Condition.}
A central difficulty of frozen-model translation is that the task gradient must flow \emph{through} the frozen predictor before reaching $\theta$.
Decomposing the gradient of Eq.~\ref{eq:composite}:
\begin{equation}
\label{eq:grad_decomp}
\nabla_\theta \mathcal{L} \;=\; \nabla_\theta \mathcal{L}_{\mathrm{task}} \;+\; \lambda_{\mathrm{fid}}\,\nabla_\theta \mathcal{L}_{\mathrm{fid}},
\end{equation}
we define the \emph{gradient alignment coefficient}:
\begin{equation}
\label{eq:alpha}
\alpha(\theta) \;=\; \cos\!\bigl(\nabla_\theta \mathcal{L}_{\mathrm{task}},\;\nabla_\theta \mathcal{L}_{\mathrm{fid}}\bigr).
\end{equation}

\begin{proposition}[Gradient Alignment Condition, informal]
\label{prop:alignment}
Stable frozen-model translation requires $\mathbb{E}[\alpha(\theta)] \geq 0$ during training.
When $\alpha < 0$, the fidelity and task gradients exhibit destructive interference: reducing $\mathcal{L}_{\mathrm{fid}}$ (moving the translation closer to identity) increases $\mathcal{L}_{\mathrm{task}}$, while reducing $\mathcal{L}_{\mathrm{task}}$ (improving predictions) increases domain divergence---violating the $\mathcal{H}\Delta\mathcal{H}$-divergence bound of \citet{bendavid2010}.
\end{proposition}

We verify this empirically across five clinical tasks.
Table~\ref{tab:alpha} reports gradient diagnostics measured at the translator parameters during the first epoch of training.\footnote{The full paper will extend these diagnostics across all three translator architectures and additional hyperparameter configurations, including per-epoch trajectories of $\alpha(\theta)$.}

\begin{table}[h]
\centering
\small
\caption{Gradient alignment diagnostics across tasks. $\|\nabla_\theta \mathcal{L}_{\mathrm{fid}}\| / \|\nabla_\theta \mathcal{L}_{\mathrm{task}}\|$ is the fidelity-to-task gradient magnitude ratio; $\alpha$ is the cosine alignment (Eq.~\ref{eq:alpha}). Positive label rate is per-timestep for AKI/Sepsis and per-stay for Mortality.}
\label{tab:alpha}
\begin{tabular}{lcccl}
\toprule
\textbf{Task} & \textbf{Label rate} & $\|\nabla \mathcal{L}_{\mathrm{fid}}\| / \|\nabla \mathcal{L}_{\mathrm{task}}\|$ & $\alpha$ & \textbf{Outcome} \\
\midrule
Mortality & 5.5\% (per-stay) & 2.8$\times$ & $+0.84$ & $+0.048$ AUROC \\
AKI & 11.95\% (per-ts) & moderate & $>0$ & $+0.056$ AUROC \\
Sepsis & 1.13\% (per-ts) & 5.7$\times$ & $-0.21$ & $+0.006$ AUROC (delta) \\
\bottomrule
\end{tabular}
\end{table}

Across the three classification tasks, the alignment coefficient is the strongest predictor of translation success: mortality ($\alpha = +0.84$) and AKI ($\alpha > 0$) succeed because fidelity acts as a step-size regulariser along the task gradient direction, whereas sepsis ($\alpha = -0.21$) suffers destructive interference.
Removing fidelity entirely ($\lambda_{\mathrm{fid}} = 0$) causes catastrophic divergence ($-0.101$ AUROC, from $0.719$ to $0.618$), confirming that the task gradient alone is too noisy to guide the translator---the fidelity term is necessary but, when misaligned, actively harmful.
This tension motivates the progression from delta-based translation (Section~\ref{sec:delta}) through shared-latent alignment (Section~\ref{sec:sl}) to retrieval-guided translation (Section~\ref{sec:retrieval}), where the retrieval mechanism provides an alternative information pathway that sidesteps the gradient alignment problem entirely.

% ============================================================
\section{Method}
\label{sec:method}
% ============================================================

We present three translator architectures of increasing capability. Each addresses a specific limitation of its predecessor, producing a clear ablation path: delta (limited by gradient alignment) $\to$ shared latent (limited by reconstruction bottleneck) $\to$ retrieval-guided (resolves both).

\paragraph{4.1\quad Delta Translator.}
\label{sec:delta}
The delta translator produces additive corrections:
\begin{equation}
\label{eq:delta}
g_\theta^{\delta}(x) \;=\; x + \Delta_\theta(x),
\end{equation}
where $\Delta_\theta$ is parameterised by a stack of \emph{AxialBlocks}---factored attention that alternates variable-wise attention (each timestep's $F$ features attend to one another) and temporal attention (each feature's $T$ timesteps attend across time).
This factorisation reduces the attention cost from $\mathcal{O}((TF)^2)$ to $\mathcal{O}(T \cdot F^2 + F \cdot T^2)$ while preserving full expressivity~\cite{ho2019axial}.
Static demographics (age, sex, height, weight) condition each block via FiLM layers~\cite{perez2018film}: $h \mapsto \gamma \odot h + \beta$ with $(\gamma, \beta)$ predicted from the static features.
Temporal attention is \emph{causal} (masked to prevent future information leakage) for per-timestep tasks (AKI, Sepsis, LoS) and \emph{bidirectional} for per-stay tasks (Mortality, KidneyFunction).

The residual formulation ensures $g_\theta$ begins near identity---early training produces small perturbations, avoiding the divergence observed with $\lambda_{\mathrm{fid}} = 0$.
However, the delta translator is fundamentally limited by Proposition~\ref{prop:alignment}: when $\alpha < 0$ (sepsis), the task and fidelity gradients create destructive interference that caps performance at $+0.006$ AUROC despite extensive hyperparameter search (13 approaches $\times$ 2 tasks; detailed in the full paper).

\paragraph{4.2\quad Shared Latent Translator.}
\label{sec:sl}
To address gradient misalignment, the shared latent (SL) translator replaces the direct input-to-input mapping with an encoder--decoder architecture that passes through a shared bottleneck $z \in \mathbb{R}^{T \times d_{\mathrm{latent}}}$.
Training proceeds in two phases.
\textbf{Phase~1} (autoencoder pretrain): the encoder--decoder is trained on MIMIC data only, with reconstruction and optional label prediction losses.
\textbf{Phase~2} (translation): the full loss combines task, reconstruction, distributional alignment, and label prediction:
\begin{equation}
\label{eq:sl_loss}
\mathcal{L}_{\mathrm{SL}} = \mathcal{L}_{\mathrm{task}} + \lambda_{\mathrm{recon}}\,\mathcal{L}_{\mathrm{recon}} + \lambda_{\mathrm{align}}\,\mathrm{MMD}(z_S, z_T) + \lambda_{\mathrm{label}}\,\mathcal{L}_{\mathrm{label\_pred}},
\end{equation}
where $\mathrm{MMD}$ is the maximum mean discrepancy with a multi-scale Gaussian kernel~\cite{gretton2012kernel} and $\mathcal{L}_{\mathrm{label\_pred}}$ provides a direct label supervision signal that bypasses the frozen predictor.
The MMD term provides distributional alignment at the latent level, replacing the input-space fidelity whose gradient opposes the task signal.

SL succeeds on tasks with sufficient label density: $+0.048$ AUROC ($+0.055$ AUCPR) on mortality (5.5\% per-stay labels), $+0.054$ AUROC ($+0.133$ AUCPR) on AKI (11.95\% per-timestep labels).
However, it fails on sepsis ($-0.017$ to $-0.043$ AUROC), exposing a \emph{reconstruction bottleneck}: the encoder--decoder must reconstruct all $T \times F$ features (e.g., $169 \times 48$) through the latent space.
At sepsis's 1.13\% positive rate, each training sequence contains ${\sim}1$ informative task gradient versus ${\sim}8{,}000$ reconstruction gradients---reconstruction dominates and the task signal cannot steer the latent representations.
Crucially, this is not a per-timestep or causality issue: AKI is also per-timestep and causal, yet SL succeeds because its 11.95\% label density provides $10\times$ more task gradient per sequence.
This controlled comparison isolates \emph{label density} as the root cause and motivates an architecture that does not rely on dense reconstruction to produce accurate translations.


\paragraph{4.3\quad Retrieval-Guided Translator.}
\label{sec:retrieval}
The retrieval-guided translator is our primary contribution.
Instead of global distributional alignment (SL's MMD) or direct input perturbation (delta), it performs \emph{instance-level retrieval}: each source timestep finds its nearest target-domain neighbours and selectively borrows feature-level information via cross-attention.
This opens a different information pathway: the model need not ``invent'' correct target-domain feature values from noisy task gradients, but can instead attend to actual target-domain exemplars.

\vspace{0.3em}
\noindent\textit{Memory Bank.}\quad
All target-domain (MIMIC) training data is encoded through the shared encoder and segmented into temporal windows of size $W$ (default $W{=}6$) with configurable stride $s$ for overlapping coverage.
Each window is mean-pooled to produce a compact descriptor:
\begin{equation}
\label{eq:memory}
M_j \;=\; \frac{1}{|W_j|} \sum_{t \in W_j} \mathrm{Enc}(x_T^{(i)})_t \;\in\; \mathbb{R}^{d_{\mathrm{latent}}},
\end{equation}
yielding a memory bank $M \in \mathbb{R}^{N_{\mathrm{windows}} \times d_{\mathrm{latent}}}$ that resides on GPU for fast retrieval.
The memory bank is rebuilt every $R$ epochs (default $R{=}3$), since the encoder evolves during training and a stale bank degrades retrieval quality.
Overlapping windows (stride $s{=}3$, yielding ${\sim}2\times$ more windows) improve performance on AKI ($+0.003$ AUROC, $+0.010$ AUCPR), but excessive overlap (stride $s{=}1$, ${\sim}2.4$M windows) degrades results ($-0.002$ AUROC), suggesting a coverage--redundancy tradeoff where moderate overlap improves recall of relevant exemplars but excessive redundancy dilutes the nearest-neighbour signal.

\vspace{0.3em}
\noindent\textit{Per-Timestep k-NN Retrieval.}\quad
For each source timestep $t$, we form a backward-looking query by mean-pooling source latents from $[\max(0, t{-}W{+}1),\; t{+}1)$, ensuring causal consistency.
Learned importance weights $\sigma(\mathbf{w}) \in [0,1]^{d_{\mathrm{latent}}}$ (sigmoid-gated, per latent dimension) modulate the distance metric:
\begin{equation}
\label{eq:retrieval_dist}
d(q_t, M_j) \;=\; \bigl\|\sqrt{\sigma(\mathbf{w})} \odot (q_t - M_j)\bigr\|_2,
\end{equation}
and the $K$ nearest windows per timestep are retrieved (default $K{=}16$).
Full timestep-level latents from each retrieved window are gathered to form the context tensor $C \in \mathbb{R}^{B \times T \times KW \times d_{\mathrm{latent}}}$.
Retrieval is per-timestep rather than per-sequence: we find that local temporal matching (window $W{=}6$) significantly outperforms wider windows ($W{=}12$: $+0.046$ vs.\ $+0.056$ AUROC on AKI), consistent with clinical time series where local physiological states are more informative than global trajectory similarity.

\vspace{0.3em}
\noindent\textit{Cross-Attention Fusion.}\quad
Retrieved context is integrated through $L$ stacked cross-attention blocks.
Each block performs three stages:
\begin{equation}
\label{eq:cross_attn}
\begin{aligned}
h_t' &= \mathrm{CrossAttn}(Q{=}h_t,\; K{=}C_t,\; V{=}C_t) + h_t, \\
\tilde{h} &= \mathrm{CausalSelfAttn}(h') + h', \\
\hat{h} &= \mathrm{FFN}(\tilde{h}) + \tilde{h},
\end{aligned}
\end{equation}
where $h_t$ is the source latent at timestep $t$, $C_t$ contains the $KW$ retrieved target latents for that timestep, and all sub-layers use pre-LayerNorm residual connections.
Stage~1 (per-timestep cross-attention) performs selective, feature-level borrowing from target exemplars---critically, this is \emph{learned}, not interpolated.
Compared to the $k$NN-MT approach of \citet{khandelwal2021knnmt}, which interpolates uniformly across all dimensions, cross-attention allows the model to attend selectively to different features from different neighbours.
Stage~2 (global causal self-attention) enforces temporal coherence across the translated sequence.

The number of cross-attention blocks $L$ is task-dependent and interacts with label density.
On AKI (11.95\% positive rate), increasing from $L{=}2$ to $L{=}3$ yields $+0.009$ AUROC---the additional capacity integrates richer target context from the dense supervision signal.
On sepsis (1.13\% positive rate), $L{=}3$ \emph{hurts} by $-0.006$ AUROC: extra parameters overfit when gradient signal is sparse.
This interaction provides a clean ablation of how architectural capacity must be matched to task supervision density.

\vspace{0.3em}
\noindent\textit{Why Retrieval Resolves Gradient Alignment.}\quad
The fundamental limitation of delta and SL translators is that the task gradient must specify both \emph{which features} to correct and \emph{what values} to produce---all by backpropagating through the frozen model's opaque hidden states.
Retrieval decouples these two problems by introducing an orthogonal information channel.
The memory bank supplies actual target-domain feature values; the model need only learn a soft \emph{selection} over them via cross-attention weights.
This constrains the output space to weighted combinations of real target exemplars, yielding a better-conditioned optimisation problem: the task gradient specifies \emph{which} neighbours are useful (a selection problem over $K$ candidates), while the neighbours themselves supply the \emph{content} (correct feature values from the target distribution).
Because the fidelity objective is now computed against outputs that are already grounded in the target domain, its gradient no longer opposes the task gradient---both objectives benefit from selecting more relevant target exemplars.

A critical design choice is detaching the source latent before querying the memory bank: $q_t = \mathrm{Enc}(x_S)_t.\mathrm{detach}()$.
Since $k$-NN retrieval involves a non-differentiable $\operatorname{argmin}$, gradients cannot propagate through the retrieval step; detaching prevents spurious gradient signals from the bank lookup.
The encoder is still updated through the downstream cross-attention and decoding losses.

The empirical results confirm this analysis.
On sepsis, where $\alpha = -0.21$ renders delta translation ineffective ($+0.006$ AUROC) and SL collapses ($-0.017$ to $-0.043$), retrieval achieves $+0.051$ AUROC---compared to $+0.006$ for the best delta configuration and negative results for the shared latent translator, establishing retrieval as the only architecture that produces clinically meaningful improvement on this task.
On AKI, retrieval ($+0.056$ AUROC, $+0.161$ AUCPR) surpasses both delta and SL, reaching an absolute AUROC of $91.14$ and AUCPR of $72.9$---above even the in-domain MIMIC-native LSTM ($89.7$ AUROC, $66.5$ AUCPR)~\cite{vandewater2024yaib}.
The AUCPR improvement ($+16.1$ points) is nearly $3\times$ larger than the AUROC gain ($+5.6$), indicating that the translator disproportionately improves precision at high recall---the regime most relevant for clinical deployment where false alarms carry real cost.

\vspace{0.3em}
\noindent\textit{FeatureGate.}\quad
We introduce a \emph{FeatureGate} module, applicable across all three paradigms: learnable per-feature sigmoid weights $\sigma(w_f) \in [0,1]$ that modulate the reconstruction loss:
\begin{equation}
\label{eq:feature_gate}
\mathcal{L}_{\mathrm{recon}}^{\mathrm{gated}} = \sum_{f=1}^{F} \sigma(w_f) \cdot \bigl(x_{\mathrm{out}}^f - x_{\mathrm{target}}^f\bigr)^2.
\end{equation}
This automatically down-weights reconstruction error on features irrelevant to the downstream task (e.g., features the frozen model ignores), focusing capacity on task-critical dimensions.
FeatureGate provides consistent gains: on mortality, SL+FG achieves $+0.048$ AUROC ($+0.055$ AUCPR) vs.\ plain SL $+0.044$ AUROC ($+0.048$ AUCPR).

\paragraph{4.4\quad Training Procedure.}
\label{sec:training}
Training follows a two-phase protocol common to SL and retrieval translators.
\textbf{Phase~1} (autoencoder pretrain, 15 epochs on MIMIC): the encoder--decoder is trained with reconstruction and label prediction losses.
For retrieval, we introduce \emph{self-retrieval pretraining}: during Phase~1, a memory bank
is built from the encoder's own MIMIC representations and refreshed periodically.
This provides non-degenerate context to cross-attention blocks from the start of training;
without self-retrieval, cross-attention receives zero tensors during Phase~1, collapsing
to a trivial pass-through that wastes the pretraining phase for the cross-attention parameters.
\textbf{Phase~2} (retrieval-guided translation, 50 epochs): the full loss is:
\begin{equation}
\label{eq:full_loss}
\mathcal{L} = \mathcal{L}_{\mathrm{task}} + \lambda_{\mathrm{recon}}\,\mathcal{L}_{\mathrm{recon}} + \lambda_{\mathrm{range}}\,\mathcal{L}_{\mathrm{range}} + \lambda_{\mathrm{align}}\,\mathrm{MMD} + \lambda_{\mathrm{target}}\,\mathcal{L}_{\mathrm{task}}^{\mathrm{MIMIC}} + \lambda_{\mathrm{label}}\,\mathcal{L}_{\mathrm{label\_pred}},
\end{equation}
where $\mathcal{L}_{\mathrm{task}}^{\mathrm{MIMIC}}$ applies the same task loss to MIMIC data through the frozen predictor, preventing representational drift.
The memory bank is rebuilt every 3 epochs; learning rate follows cosine annealing~\cite{loshchilov2017sgdr}; gradient norms are clipped to $1.0$; early stopping with patience 15 monitors the primary metric.

All data follows the YAIB preprocessing pipeline~\cite{vandewater2024yaib}: standardized variable selection (48 physiological features common to both databases), forward-fill imputation with missing indicators, per-feature normalization, and consistent patient-level train/validation/test splits.
This ensures that distributional differences between eICU and MIMIC-IV reflect institutional variation rather than preprocessing artifacts.
On top of YAIB's per-database normalization, we apply \emph{cross-domain normalisation}: an affine transform that matches source per-feature mean and variance to target statistics prior to translation.
This is computed once from training data and stored in the checkpoint, ensuring consistent normalisation at inference.

Phase~1 checkpoints are reusable: because pretraining depends only on the target data and model architecture (not Phase~2 hyperparameters), a single pretrain checkpoint serves all downstream experiments for a given task.
We verify this empirically: across 9 AKI configurations varying $K$, $W$, $L$, and loss weights, the Phase~1 label prediction loss converges to $0.462$--$0.466$, confirming Phase~2 independence.
